\documentclass[10pt,a4paper]{amsart}
\usepackage{amsmath,amssymb,amsthm,graphicx}
% Retained minimal preamble for the Stacks Project chapter more-algebra.tex (More on Algebra).
% The chapter file pinned-original-source.tex carries only \input{preamble} on its first line, so no
% preamble text lives inside the chapter itself.  The Stacks Project uses one single shared preamble for
% every chapter; the snapshot stages it identically as sources/stacks-chow/preamble.tex and
% sources/stacks-etale-cohomology/preamble.tex.  The block below is copied mechanically, line for line,
% from sources/stacks-chow/preamble.tex lines 14-18 (the Xy-pic setup) and 154-217 (the native theorem
% declarations together with the mathematical macros).  Nothing was added, removed or redefined; the
% wrapper class and the amsmath/amssymb/amsthm/graphicx base packages are supplied by the assembler.
% Chapter: more-algebra.tex -- More on Algebra.
% For commutative diagrams we use Xy-pic
\usepackage[all]{xy}

% We use 2cell for 2-commutative diagrams.
\xyoption{2cell}

% Theorem environments.
%
\theoremstyle{plain}
\newtheorem{theorem}[subsection]{Theorem}
\newtheorem{proposition}[subsection]{Proposition}
\newtheorem{lemma}[subsection]{Lemma}

\theoremstyle{definition}
\newtheorem{definition}[subsection]{Definition}
\newtheorem{example}[subsection]{Example}
\newtheorem{exercise}[subsection]{Exercise}
\newtheorem{situation}[subsection]{Situation}

\theoremstyle{remark}
\newtheorem{remark}[subsection]{Remark}
\newtheorem{remarks}[subsection]{Remarks}

\numberwithin{equation}{subsection}

% Macros
%
\def\lim{\mathop{\mathrm{lim}}\nolimits}
\def\colim{\mathop{\mathrm{colim}}\nolimits}
\def\Spec{\mathop{\mathrm{Spec}}}
\def\Hom{\mathop{\mathrm{Hom}}\nolimits}
\def\Ext{\mathop{\mathrm{Ext}}\nolimits}
\def\SheafHom{\mathop{\mathcal{H}\!\mathit{om}}\nolimits}
\def\SheafExt{\mathop{\mathcal{E}\!\mathit{xt}}\nolimits}
\def\Sch{\mathit{Sch}}
\def\Mor{\mathop{\mathrm{Mor}}\nolimits}
\def\Ob{\mathop{\mathrm{Ob}}\nolimits}
\def\Sh{\mathop{\mathit{Sh}}\nolimits}
\def\NL{\mathop{N\!L}\nolimits}
\def\CH{\mathop{\mathrm{CH}}\nolimits}
\def\proetale{{pro\text{-}\acute{e}tale}}
\def\etale{{\acute{e}tale}}
\def\QCoh{\mathit{QCoh}}
\def\Ker{\mathop{\mathrm{Ker}}}
\def\Im{\mathop{\mathrm{Im}}}
\def\Coker{\mathop{\mathrm{Coker}}}
\def\Coim{\mathop{\mathrm{Coim}}}

% Boxtimes
%
\DeclareMathSymbol{\boxtimes}{\mathbin}{AMSa}{"02}

%
% Macros for moduli stacks/spaces
%
\def\QCohstack{\mathcal{QC}\!\mathit{oh}}
\def\Cohstack{\mathcal{C}\!\mathit{oh}}
\def\Spacesstack{\mathcal{S}\!\mathit{paces}}
\def\Quotfunctor{\mathrm{Quot}}
\def\Hilbfunctor{\mathrm{Hilb}}
\def\Curvesstack{\mathcal{C}\!\mathit{urves}}
\def\Polarizedstack{\mathcal{P}\!\mathit{olarized}}
\def\Complexesstack{\mathcal{C}\!\mathit{omplexes}}
% \Pic is the operator that assigns to X its picard group, usage \Pic(X)
% \Picardstack_{X/B} denotes the Picard stack of X over B
% \Picardfunctor_{X/B} denotes the Picard functor of X over B
\def\Pic{\mathop{\mathrm{Pic}}\nolimits}
\def\Picardstack{\mathcal{P}\!\mathit{ic}}
\def\Picardfunctor{\mathrm{Pic}}
\def\Deformationcategory{\mathcal{D}\!\mathit{ef}}

% --- appended 2026-10-12 (Chief): the project own definitions of the reference/slogan/history
% display environments, copied mechanically from the shared Stacks preamble lines 8-12
% (\usepackage{verbatim} plus the three \newenvironment lines, which the chapter preamble keeps
% verbatim from the project preamble). Nothing added, removed or redefined: these environments
% carry no mathematical content and the project itself typesets them as comments. Without these
% three lines a counted statement whose own text uses them fails with Environment slogan undefined.
\usepackage{verbatim}
\newenvironment{reference}{\comment}{\endcomment}
\newenvironment{slogan}{\comment}{\endcomment}
\newenvironment{history}{\comment}{\endcomment}
\usepackage{cleveref}
\crefname{originallemma}{original lemma}{original lemmas}
\Crefname{originallemma}{Original lemma}{Original lemmas}
\makeatletter
\newlabel{more-algebra-lemma-power-series-ring-subfields}{{07P5}{source}{}{sourcepointer-1}{}}
\newlabel{more-algebra-lemma-power-series-ring-subfields@cref}{{[originallemma][0][]07P5}{[1][0][]source}{}{}{}}
\makeatother
\crefname{originallemma}{original lemma}{original lemmas}
\Crefname{originallemma}{Original lemma}{Original lemmas}
\makeatletter
\newlabel{lemma-intersection-subfields}{{07P4}{source}{}{sourcepointer-2}{}}
\newlabel{lemma-intersection-subfields@cref}{{[originallemma][0][]07P4}{[1][0][]source}{}{}{}}
\makeatother
\crefname{originallemma}{original lemma}{original lemmas}
\Crefname{originallemma}{Original lemma}{Original lemmas}
\makeatletter
\newlabel{algebra-lemma-finite-after-completion}{{0394}{source}{}{sourcepointer-3}{}}
\newlabel{algebra-lemma-finite-after-completion@cref}{{[originallemma][0][]0394}{[1][0][]source}{}{}{}}
\makeatother
\title{A filtering family of subfields of a power series fraction field with intersection the p-th powers}
\author{The Stacks Project Authors}
\date{stacks-project:more-algebra.tex:local-snapshot-20261012; GNU FDL 1.2 or later}
\begin{document}
\maketitle
\setcounter{subsection}{1}
\noindent\textit{Editorial note.} All delivered mathematical text is the author's own, copied from the pinned Stacks Project chapter file more-algebra.tex; no statement or proof step was rewritten, reordered or removed.\par\medskip
\noindent\textit{Editorial note.} The chapter file pinned-original-source.tex carries only \textbackslash{}input\{preamble\} on its first line, so no preamble text lives inside the chapter. The retained file sources/stacks-more-algebra/retained-preamble.tex is a mechanical, line-for-line copy of the project's single shared preamble as staged in the snapshot (Xy-pic setup plus theorem declarations and mathematical macros). Nothing was added, removed or redefined, and the excerpt serves only as preamble, never as a span coordinate.\par\medskip
\noindent\textit{Editorial note.} Both counted spans start at the \textbackslash{}begin line of their environment, so the statement and the author's proof are delivered whole; surrounding prose of the section is deliberately outside the counted spans and is carried as standard prerequisites instead.\par\medskip
\noindent\textit{Editorial note.} The citation fields (source\_id, source\_version, authors, source\_url, license.name, source\_title) are identical to those of every other card this lane has produced from this chapter, so the batch selector treats them as one source.\par\medskip
\noindent\textit{Allowed standard prerequisites (editorial).}\begin{itemize}
\item p-basis of a field of characteristic p. the statement counted by the sibling card S1-algebra-p-basis at 11362-11415. Supplies the indexing set and the description of k\_J.
\item Finiteness of a power series extension after a finite extension of the coefficient field. the author cites Algebra, Lemma algebra-lemma-finite-after-completion at proof lines 11537-11539. Gives finiteness of the ring extensions A\_J to A.
\item The intersection statement for filtering families of subfields, whose three hypotheses are verified here. the statement counted by the sibling card S1-algebra-intersection-subfields-subspace at 11417-11455, cited at proof lines 11542-11544. Supplies the list of properties to be checked and the use to which the family is put elsewhere.
\end{itemize}\medskip
\ifdefined\hypertarget\hypertarget{sourcepointer-1}{}\fi
\noindent\textit{Source reference.} Original lemma 07P5: focused-seeds/AGP/sources/stacks-snapshot-20261011/tags.txt, source lines 9785--9785. The lemma's own canonical tag, quoted so that readers can resolve the label inside the delivered text; it points at the counted statement itself and carries no proof content.\par\medskip
\ifdefined\hypertarget\hypertarget{sourcepointer-2}{}\fi
\noindent\textit{Source reference.} Original lemma 07P4: focused-seeds/AGP/sources/stacks-snapshot-20261011/tags.txt, source lines 9784--9784. Navigation pointer to a statement the counted proof refers to but does not contain; nothing load-bearing is delegated to it.\par\medskip
\ifdefined\hypertarget\hypertarget{sourcepointer-3}{}\fi
\noindent\textit{Source reference.} Original lemma 0394: focused-seeds/AGP/sources/stacks-snapshot-20261011/tags.txt, source lines 4259--4259. Navigation pointer to a statement the counted proof refers to but does not contain; nothing load-bearing is delegated to it.\par\medskip
\begin{lemma}
\label{lemma-power-series-ring-subfields}
Let $k$ be a field of characteristic $p > 0$. Let $\{x_i\}_{i \in I}$
be a $p$-basis for $k$. Let $n, m \geq 0$.
Let $K$ be the fraction field of $A = k[[x_1, \ldots, x_n]][y_1, \ldots, y_m]$.
Let $J$ be a finite subset of $I$. Consider the subfield
$k/k_J/k^p$ generated by $k^p$ and $x_i$ with $i \in I \setminus J$.
The fraction fields $K_J$ of
$$
A_J = k_J[[x_1^p, \ldots, x_n^p]][y_1^p, \ldots, y_m^p]
$$
form a family of subfields of $K$ as in
Lemma \ref{lemma-intersection-subfields}. Moreover, each of the ring extensions
$A_J \subset A$ is finite.
\end{lemma}

\begin{proof}
Since $k/k_J$ is finite, the ring extension
$k_J[[x_1^p, \ldots, x_d^p]] \subset k[[x_1, \ldots, x_d]]$ is finite
by Algebra, Lemma \ref{algebra-lemma-finite-after-completion}.
This implies that $A_J \to A$ is finite.

\medskip\noindent
Let us check properties (1), (2), (3) of
Lemma \ref{lemma-intersection-subfields}.
Proof of (1). For $a \in A$ we see that $a^p \in A_J$.
Hence $K^p \subset K_J$. Proof of (2).
Suppose that $f/g^p \in K$, $f, g \in A$, $g \not = 0$
is contained in $K_J$ for every choice of $J$.
Fix $J$ for the moment. Since $f/g^p \in K_J$ we can write
$f/g^p = a/b^p$ with $a \in A_J$ and $b \in A$ nonzero.
Hence $b^p f \in A_J$. For any $A_J$-derivation $D : A \to A$ we see
that $0 = D(b^pf) = b^p D(f)$ hence $D(f) = 0$ as $A$ is a domain.
Taking $D = \partial_{x_i}$ and $D = \partial_{y_j}$ we conclude
that $f \in k[[x_1^p, \ldots, x_n^p]][y_1^p, \ldots, y_d^p]$.
Applying a $k_J$-derivation $\theta : k \to k$
we similarly conclude that all coefficients of $f$
are in $k_J$, i.e., $f \in A_J$. Since it is clear that
$A^p = \bigcap\nolimits_J A_J$ where $J$ ranges over all subfields
as in the lemma we conclude $f \in A^p$ as desired.
Proof of (3). This is clear because $K_{J \cup J'} \subset K_J \cap K_{J'}$.
\end{proof}

\end{document}
